%%%%%%%%%%%%%%%%%%%%%%%%%%%%%%%%%%%%%%%%%%%%%%%%%%%%%%%%%%%%%%%%%%%%%%
% ALORA
% Automatic Learning for Oral Resource-poor African Languages
%
% Note de cadrage scientifique
% Version 0.1
%%%%%%%%%%%%%%%%%%%%%%%%%%%%%%%%%%%%%%%%%%%%%%%%%%%%%%%%%%%%%%%%%%%%%%

\documentclass[11pt,a4paper]{article}

%%%%%%%%%%%%%%%%%%%%%%%%%%%%%%%%%%%%%%%%%%%%%%%%%%%%%%%%%%%%%%%%%%%%%%
%% Encodage
%%%%%%%%%%%%%%%%%%%%%%%%%%%%%%%%%%%%%%%%%%%%%%%%%%%%%%%%%%%%%%%%%%%%%%

\usepackage[T1]{fontenc}
\usepackage[utf8]{inputenc}
\usepackage[french]{babel}

%%%%%%%%%%%%%%%%%%%%%%%%%%%%%%%%%%%%%%%%%%%%%%%%%%%%%%%%%%%%%%%%%%%%%%
%% Mise en page
%%%%%%%%%%%%%%%%%%%%%%%%%%%%%%%%%%%%%%%%%%%%%%%%%%%%%%%%%%%%%%%%%%%%%%

\usepackage[a4paper,margin=2.5cm]{geometry}

%%%%%%%%%%%%%%%%%%%%%%%%%%%%%%%%%%%%%%%%%%%%%%%%%%%%%%%%%%%%%%%%%%%%%%
%% Packages
%%%%%%%%%%%%%%%%%%%%%%%%%%%%%%%%%%%%%%%%%%%%%%%%%%%%%%%%%%%%%%%%%%%%%%

\usepackage{graphicx}
\usepackage{float}
\usepackage{booktabs}
\usepackage{longtable}
\usepackage{array}
\usepackage{enumitem}
\usepackage{setspace}
\usepackage{hyperref}
\usepackage{xcolor}
\usepackage{titlesec}
\usepackage{tikz}
\usepackage{tcolorbox}
\usepackage{amsmath}
\usepackage{amssymb}

%%%%%%%%%%%%%%%%%%%%%%%%%%%%%%%%%%%%%%%%%%%%%%%%%%%%%%%%%%%%%%%%%%%%%%
%% Bibliographie
%%%%%%%%%%%%%%%%%%%%%%%%%%%%%%%%%%%%%%%%%%%%%%%%%%%%%%%%%%%%%%%%%%%%%%

\usepackage[
backend=biber,
style=ieee
]{biblatex}

\addbibresource{references.bib}

%%%%%%%%%%%%%%%%%%%%%%%%%%%%%%%%%%%%%%%%%%%%%%%%%%%%%%%%%%%%%%%%%%%%%%
%% Hyperliens
%%%%%%%%%%%%%%%%%%%%%%%%%%%%%%%%%%%%%%%%%%%%%%%%%%%%%%%%%%%%%%%%%%%%%%

\hypersetup{
colorlinks=true,
linkcolor=blue,
citecolor=blue,
urlcolor=blue
}

%%%%%%%%%%%%%%%%%%%%%%%%%%%%%%%%%%%%%%%%%%%%%%%%%%%%%%%%%%%%%%%%%%%%%%
%% Espacement
%%%%%%%%%%%%%%%%%%%%%%%%%%%%%%%%%%%%%%%%%%%%%%%%%%%%%%%%%%%%%%%%%%%%%%

\setstretch{1.2}

\setlength{\parindent}{0pt}
\setlength{\parskip}{0.6em}

%%%%%%%%%%%%%%%%%%%%%%%%%%%%%%%%%%%%%%%%%%%%%%%%%%%%%%%%%%%%%%%%%%%%%%
%% Titres
%%%%%%%%%%%%%%%%%%%%%%%%%%%%%%%%%%%%%%%%%%%%%%%%%%%%%%%%%%%%%%%%%%%%%%

\titleformat{\section}
{\Large\bfseries}
{\thesection}
{1em}
{}

\titleformat{\subsection}
{\large\bfseries}
{\thesubsection}
{1em}
{}

%%%%%%%%%%%%%%%%%%%%%%%%%%%%%%%%%%%%%%%%%%%%%%%%%%%%%%%%%%%%%%%%%%%%%%
%% Encadrés
%%%%%%%%%%%%%%%%%%%%%%%%%%%%%%%%%%%%%%%%%%%%%%%%%%%%%%%%%%%%%%%%%%%%%%

\newtcolorbox{objectif}{
colback=blue!5,
colframe=blue!70!black,
title=Objectif
}

\newtcolorbox{discussion}{
colback=green!5,
colframe=green!60!black,
title=Point de discussion
}

%%%%%%%%%%%%%%%%%%%%%%%%%%%%%%%%%%%%%%%%%%%%%%%%%%%%%%%%%%%%%%%%%%%%%%
%%%%%%%%%%%%%%%%%%%%%%%%%%%%%%%%%%%%%%%%%%%%%%%%%%%%%%%%%%%%%%%%%%%%%%

\begin{document}

%%%%%%%%%%%%%%%%%%%%%%%%%%%%%%%%%%%%%%%%%%%%%%%%%%%%%%%%%%%%%%%%%%%%%%
%% Couverture
%%%%%%%%%%%%%%%%%%%%%%%%%%%%%%%%%%%%%%%%%%%%%%%%%%%%%%%%%%%%%%%%%%%%%%

\begin{titlepage}

\centering

\vspace*{1.8cm}

{\Huge\bfseries ALORA}

\vspace{0.5cm}

{\Large
Automatic Learning for Oral Resource-poor African Languages
}

\vspace{0.5cm}

{\large
Apprentissage automatique des langues orales africaines peu dotées
}

\vspace{2cm}

{\LARGE\bfseries
Note de cadrage scientifique
}

\vspace{0.3cm}

Version 0.1

\vfill

\textbf{Projet présenté par}

\vspace{0.3cm}

{\Large
Mamadou Tahirou Diallo
}

Étudiant ingénieur

Institut National des Postes et Télécommunications (INPT)

\vfill

\textbf{Objet du document}

\vspace{0.2cm}

Proposition de collaboration scientifique

Travail de recherche — Été 2026

\vfill

\today

\end{titlepage}

%%%%%%%%%%%%%%%%%%%%%%%%%%%%%%%%%%%%%%%%%%%%%%%%%%%%%%%%%%%%%%%%%%%%%%
%% Sommaire
%%%%%%%%%%%%%%%%%%%%%%%%%%%%%%%%%%%%%%%%%%%%%%%%%%%%%%%%%%%%%%%%%%%%%%

\tableofcontents

\newpage

%%%%%%%%%%%%%%%%%%%%%%%%%%%%%%%%%%%%%%%%%%%%%%%%%%%%%%%%%%%%%%%%%%%%%%
%% Résumé
%%%%%%%%%%%%%%%%%%%%%%%%%%%%%%%%%%%%%%%%%%%%%%%%%%%%%%%%%%%%%%%%%%%%%%

\section*{Résumé}

\addcontentsline{toc}{section}{Résumé}

Le présent document expose la vision initiale du projet
\textbf{ALORA (Automatic Learning for Oral Resource-poor African Languages)}.

L'objectif général est d'étudier comment les méthodes modernes
d'intelligence artificielle peuvent apprendre progressivement une langue
essentiellement orale à partir d'un corpus de parole, découvrir sa
structure phonologique, puis assister à la construction d'une
représentation écrite cohérente.

Le soussou constitue le premier cas d'étude retenu.

Ce document n'a pas pour objectif de proposer une méthodologie
définitive, mais de présenter une première réflexion scientifique,
d'identifier les principaux défis du projet et de solliciter les conseils
d'un enseignant-chercheur afin d'orienter les choix méthodologiques avant
le début des développements.

\newpage

%%%%%%%%%%%%%%%%%%%%%%%%%%%%%%%%%%%%%%%%%%%%%%%%%%%%%%%%%%%%%%%%%%%%%%
%% Sections
%%%%%%%%%%%%%%%%%%%%%%%%%%%%%%%%%%%%%%%%%%%%%%%%%%%%%%%%%%%%%%%%%%%%%%

%%%%%%%%%%%%%%%%%%%%%%%%%%%%%%%%%%%%%%%%%%%%%%%%%%%%%%%%%%%%%%%%%%%%%%
\section{Contexte et motivation}
%%%%%%%%%%%%%%%%%%%%%%%%%%%%%%%%%%%%%%%%%%%%%%%%%%%%%%%%%%%%%%%%%%%%%%

Depuis plusieurs années, les progrès réalisés en intelligence artificielle
ont profondément transformé le traitement automatique de la parole et des
langues. Des modèles tels que Whisper, wav2vec~2.0, HuBERT ou encore MMS
(Massively Multilingual Speech) ont démontré qu'il était désormais possible
d'apprendre des représentations linguistiques riches directement à partir
de grands volumes de données audio.

Cependant, ces avancées concernent principalement des langues disposant de
ressources importantes : corpus annotés, dictionnaires, systèmes
d'écriture standardisés et importantes quantités de données numériques.

À l'inverse, de nombreuses langues africaines restent encore très peu
représentées dans les technologies actuelles. Elles sont souvent qualifiées
de \emph{langues peu dotées} (\emph{low-resource languages}), car elles ne
disposent que de très peu de ressources exploitables pour entraîner des
modèles d'intelligence artificielle.

Le soussou, parlé principalement en Guinée, constitue un exemple
particulièrement intéressant. Malgré un nombre important de locuteurs,
cette langue demeure largement sous-représentée dans les grands corpus
internationaux utilisés aujourd'hui pour le développement des modèles de
traitement de la parole.

\vspace{0.5cm}

Cette situation soulève plusieurs difficultés.

Les systèmes modernes de reconnaissance vocale ne prennent généralement
pas en charge le soussou.

Les modèles de traduction automatique sont très limités.

Les assistants vocaux actuels sont incapables de comprendre ou de produire
cette langue.

Enfin, les ressources linguistiques disponibles sont insuffisantes pour
développer rapidement de nouvelles applications.

\vspace{0.5cm}

Pourtant, les récents progrès des modèles auto-supervisés montrent qu'il
n'est plus nécessaire de disposer d'énormes quantités de données annotées
pour apprendre des représentations utiles de la parole.

Cette évolution ouvre de nouvelles perspectives pour les langues peu
dotées et conduit naturellement à la question suivante :

\begin{center}
\fbox{
\parbox{0.9\textwidth}{
\centering
\textbf{
Une intelligence artificielle peut-elle apprendre progressivement une
langue essentiellement orale à partir de données de parole collectées
auprès de locuteurs natifs, puis utiliser cette connaissance pour assister
à la création d'outils linguistiques adaptés à cette langue ?
}
}
}
\end{center}

Cette interrogation constitue le point de départ du projet
\textbf{ALORA}.

\vspace{0.5cm}

Le choix du soussou comme premier cas d'étude ne repose pas uniquement sur
le fait qu'il s'agisse de ma langue maternelle. Il répond également à une
motivation scientifique.

En effet, travailler sur une langue peu dotée permet d'étudier les limites
des méthodes actuelles d'apprentissage automatique lorsque les ressources
sont rares, tout en développant des méthodologies susceptibles d'être
réutilisées pour d'autres langues présentant des caractéristiques
similaires.

Ainsi, le soussou constitue un terrain d'expérimentation pertinent pour
explorer de nouvelles approches en traitement automatique de la parole et
en linguistique computationnelle.

\vspace{0.5cm}

Ce projet ne cherche donc pas uniquement à développer une application
destinée au soussou.

L'ambition est plus générale : proposer une méthodologie permettant de
faciliter la création de ressources numériques pour des langues orales peu
dotées, en combinant collecte collaborative de données, intelligence
artificielle, traitement de la parole et apprentissage automatique.

Cette réflexion s'inscrit dans une perspective de recherche à long terme,
dont la première étape consiste à construire des bases méthodologiques
solides avant le développement des modèles d'intelligence artificielle.

%%%%%%%%%%%%%%%%%%%%%%%%%%%%%%%%%%%%%%%%%%%%%%%%%%%%%%%%%%%%%%%%%%%%%%
\section{Vision générale du projet ALORA}
%%%%%%%%%%%%%%%%%%%%%%%%%%%%%%%%%%%%%%%%%%%%%%%%%%%%%%%%%%%%%%%%%%%%%%

Le projet ALORA (\textit{Automatic Learning for Oral Resource-poor African Languages})
a pour ambition de proposer une méthodologie générale permettant de développer
progressivement des ressources numériques pour les langues orales peu dotées.

Plutôt que de commencer par la construction d'un modèle de traduction ou d'un
assistant conversationnel, l'approche proposée consiste à reproduire, autant
que possible, les grandes étapes de l'apprentissage d'une langue.

Autrement dit, avant d'espérer comprendre ou générer une langue, un système
d'intelligence artificielle doit d'abord apprendre à écouter, observer,
identifier les régularités de la parole, construire progressivement une
représentation interne de la langue, puis exploiter cette connaissance pour
développer des outils linguistiques.

Cette vision conduit naturellement à une architecture organisée en plusieurs
phases complémentaires.

%%%%%%%%%%%%%%%%%%%%%%%%%%%%%%%%%%%%%%%%%%%%%%%%%%%%%%%%%%%%%%%%%%%%%%

\subsection{Vision globale}

\begin{figure}[H]
\centering

\begin{tikzpicture}[
node distance=1.3cm,
every node/.style={
draw,
rounded corners,
align=center,
minimum width=8cm,
minimum height=0.9cm,
fill=blue!6
},
>=latex
]

\node (collecte) {Collecte collaborative de données orales};

\node (qualite) [below of=collecte]
{Contrôle qualité et structuration du corpus};

\node (prep) [below of=qualite]
{Prétraitement des données audio};

\node (speech) [below of=prep]
{Apprentissage de représentations de la parole};

\node (phono) [below of=speech]
{Découverte automatique des structures phonologiques};

\node (langue) [below of=phono]
{Apprentissage progressif de la langue};

\node (alphabet) [below of=langue]
{Induction d'un système d'écriture};

\node (applications) [below of=alphabet]
{Applications : transcription, traduction, assistants vocaux};

\draw[->] (collecte)--(qualite);
\draw[->] (qualite)--(prep);
\draw[->] (prep)--(speech);
\draw[->] (speech)--(phono);
\draw[->] (phono)--(langue);
\draw[->] (langue)--(alphabet);
\draw[->] (alphabet)--(applications);

\end{tikzpicture}

\caption{Vision générale du projet ALORA}
\end{figure}

%%%%%%%%%%%%%%%%%%%%%%%%%%%%%%%%%%%%%%%%%%%%%%%%%%%%%%%%%%%%%%%%%%%%%%

\subsection{Principe général}

Contrairement aux approches classiques qui supposent l'existence d'une langue
déjà écrite et de nombreuses ressources annotées, ALORA cherche à construire
ces ressources progressivement.

Le projet débute donc par la création d'un corpus oral de qualité, obtenu grâce
à une collecte collaborative auprès de locuteurs natifs.

Une fois ce corpus constitué, plusieurs traitements successifs permettront de
préparer les données pour l'apprentissage automatique :

\begin{itemize}

\item nettoyage des enregistrements ;

\item normalisation des formats audio ;

\item organisation des métadonnées ;

\item contrôle qualité ;

\item préparation du corpus pour les futurs modèles.

\end{itemize}

Ces premières étapes constituent les fondations du projet. Elles conditionnent
directement la qualité des modèles développés par la suite.

%%%%%%%%%%%%%%%%%%%%%%%%%%%%%%%%%%%%%%%%%%%%%%%%%%%%%%%%%%%%%%%%%%%%%%

\subsection{Une progression par niveaux}

L'objectif n'est pas de résoudre simultanément tous les problèmes liés à une
langue peu dotée.

Au contraire, ALORA adopte une stratégie incrémentale.

Chaque étape produit des résultats exploitables qui servent de base à la
suivante.

\begin{center}

\begin{tabular}{cl}

\toprule

Niveau & Objectif principal \\

\midrule

1 & Construire un corpus oral fiable \\

2 & Préparer les données pour l'apprentissage \\

3 & Apprendre des représentations de la parole \\

4 & Identifier automatiquement les unités phonologiques \\

5 & Comprendre progressivement la structure de la langue \\

6 & Étudier la possibilité de proposer une représentation écrite \\

7 & Développer des applications linguistiques \\

\bottomrule

\end{tabular}

\end{center}

Cette progression permet de limiter les risques scientifiques tout en produisant
des résultats exploitables dès les premières phases du projet.

%%%%%%%%%%%%%%%%%%%%%%%%%%%%%%%%%%%%%%%%%%%%%%%%%%%%%%%%%%%%%%%%%%%%%%

\subsection{Positionnement du travail proposé}

La totalité de cette vision représente un projet de recherche pouvant
s'étendre sur plusieurs années.

Le travail envisagé durant l'été ne constitue donc que la première étape de
cette feuille de route.

L'objectif est de construire des bases suffisamment solides pour permettre le
développement des étapes suivantes dans le cadre du semestre S5, puis du Projet
de Fin d'Études (PFE).

À ce stade, les questions relatives à la découverte automatique des unités
phonologiques, à l'apprentissage de la langue ou encore à l'induction d'un
système d'écriture restent ouvertes. Elles feront l'objet d'une étude
progressive au fur et à mesure de l'avancement du projet.

\begin{objectif}

L'objectif principal de cette première phase n'est pas de développer un modèle
d'intelligence artificielle performant, mais de concevoir une méthodologie de
collecte et d'organisation des données suffisamment robuste pour servir de
fondation à l'ensemble du projet ALORA.

\end{objectif}

%%%%%%%%%%%%%%%%%%%%%%%%%%%%%%%%%%%%%%%%%%%%%%%%%%%%%%%%%%%%%%%%%%%%%%
\section{Travail proposé durant l'été}
%%%%%%%%%%%%%%%%%%%%%%%%%%%%%%%%%%%%%%%%%%%%%%%%%%%%%%%%%%%%%%%%%%%%%%

Le projet ALORA étant conçu comme un projet de recherche à long terme,
il apparaît nécessaire de commencer par la construction de fondations
solides avant d'aborder le développement des modèles d'intelligence
artificielle.

L'objectif de cette première phase n'est donc pas de développer un
modèle de reconnaissance vocale ou de traduction, mais de construire
une infrastructure scientifique permettant de collecter, organiser et
préparer les données nécessaires aux travaux futurs.

%%%%%%%%%%%%%%%%%%%%%%%%%%%%%%%%%%%%%%%%%%%%%%%%%%%%%%%%%%%%%%%%%%%%%%

\subsection{Objectifs de cette première phase}

Les objectifs principaux sont les suivants :

\begin{enumerate}

\item Définir une méthodologie rigoureuse de collecte de données
orales.

\item Concevoir une plateforme simple permettant la participation de
locuteurs natifs.

\item Définir la structure du corpus ainsi que les métadonnées à
conserver.

\item Concevoir un pipeline de validation et de contrôle qualité des
enregistrements.

\item Produire un premier corpus exploitable pour les expérimentations
menées durant le semestre S5.

\end{enumerate}

%%%%%%%%%%%%%%%%%%%%%%%%%%%%%%%%%%%%%%%%%%%%%%%%%%%%%%%%%%%%%%%%%%%%%%

\subsection{Plateforme de collecte}

Une plateforme web légère sera développée afin de permettre à des
locuteurs natifs d'enregistrer des mots ou des phrases en soussou.

Le principe envisagé est le suivant :

\begin{itemize}

\item présentation d'un mot ou d'une phrase en français ;

\item enregistrement de la traduction orale par le participant ;

\item possibilité de réenregistrer avant validation ;

\item enregistrement automatique des métadonnées.

\end{itemize}

Cette approche permettra de constituer progressivement un corpus
contenant plusieurs réalisations d'une même phrase par différents
locuteurs.

%%%%%%%%%%%%%%%%%%%%%%%%%%%%%%%%%%%%%%%%%%%%%%%%%%%%%%%%%%%%%%%%%%%%%%

\subsection{Métadonnées envisagées}

Chaque enregistrement devra être accompagné d'informations permettant
de faciliter les analyses futures.

À titre indicatif, les métadonnées suivantes sont envisagées :

\begin{itemize}

\item identifiant du locuteur (anonyme) ;

\item âge (optionnel) ;

\item sexe (optionnel) ;

\item région d'origine ;

\item langue maternelle ;

\item langues parlées ;

\item date d'enregistrement ;

\item qualité estimée de l'audio ;

\item texte français correspondant ;

\item durée de l'enregistrement.

\end{itemize}

Cette liste reste volontairement ouverte et pourra évoluer selon les
recommandations du laboratoire.

%%%%%%%%%%%%%%%%%%%%%%%%%%%%%%%%%%%%%%%%%%%%%%%%%%%%%%%%%%%%%%%%%%%%%%

\subsection{Organisation des données}

Une attention particulière sera portée à l'organisation du corpus afin
de faciliter son exploitation.

Les travaux prévus comprennent notamment :

\begin{itemize}

\item définition de l'arborescence des données ;

\item choix des formats audio ;

\item convention de nommage des fichiers ;

\item organisation des métadonnées ;

\item préparation d'un format facilement exploitable par les outils de
Machine Learning.

\end{itemize}

%%%%%%%%%%%%%%%%%%%%%%%%%%%%%%%%%%%%%%%%%%%%%%%%%%%%%%%%%%%%%%%%%%%%%%

\subsection{Contrôle qualité}

Avant toute utilisation scientifique, les données devront être
contrôlées afin de garantir leur qualité.

Plusieurs mécanismes sont envisagés :

\begin{itemize}

\item détection des enregistrements silencieux ;

\item contrôle de la durée ;

\item détection de fichiers corrompus ;

\item suppression des doublons ;

\item vérification de la cohérence entre les métadonnées et
l'enregistrement.

\end{itemize}

%%%%%%%%%%%%%%%%%%%%%%%%%%%%%%%%%%%%%%%%%%%%%%%%%%%%%%%%%%%%%%%%%%%%%%

\subsection{Résultats attendus}

À l'issue de cette première phase, les livrables attendus sont les
suivants :

\begin{enumerate}

\item une plateforme fonctionnelle de collecte de données ;

\item une première version du corpus oral ;

\item une documentation décrivant la structure des données ;

\item un pipeline de validation et de préparation du corpus ;

\item un ensemble de recommandations méthodologiques pour les étapes
suivantes du projet.

\end{enumerate}

%%%%%%%%%%%%%%%%%%%%%%%%%%%%%%%%%%%%%%%%%%%%%%%%%%%%%%%%%%%%%%%%%%%%%%

\begin{discussion}

Cette première phase représente principalement un travail de
construction méthodologique.

Je souhaiterais profiter de cette période pour bénéficier de votre
expérience afin de valider les choix scientifiques concernant la
collecte des données, l'organisation du corpus, les métadonnées à
conserver ainsi que les outils les plus adaptés avant le démarrage des
travaux de modélisation.

\end{discussion}

%%%%%%%%%%%%%%%%%%%%%%%%%%%%%%%%%%%%%%%%%%%%%%%%%%%%%%%%%%%%%%%%%%%%%%
\section{Questions et points de discussion}
%%%%%%%%%%%%%%%%%%%%%%%%%%%%%%%%%%%%%%%%%%%%%%%%%%%%%%%%%%%%%%%%%%%%%%

L'objectif de cette note n'est pas de présenter une méthodologie figée.
Au contraire, plusieurs choix techniques et scientifiques restent ouverts
et mériteraient d'être discutés avec un enseignant-chercheur avant le
démarrage effectif du projet.

Les questions suivantes constituent les principaux points sur lesquels
je souhaiterais bénéficier de votre expérience.

%%%%%%%%%%%%%%%%%%%%%%%%%%%%%%%%%%%%%%%%%%%%%%%%%%%%%%%%%%%%%%%%%%%%%%

\subsection{1. Constitution du corpus}

\begin{discussion}

\begin{itemize}

\item Quel volume de données serait raisonnable pour une première
expérimentation scientifique ?

\item Est-il préférable de collecter un grand nombre de locuteurs avec
peu d'enregistrements chacun, ou un nombre plus réduit de locuteurs avec
davantage d'enregistrements ?

\item Faut-il commencer par des mots isolés, des phrases simples ou
directement des conversations spontanées ?

\item Existe-t-il des méthodologies de collecte que vous recommanderiez
pour une langue peu dotée comme le soussou ?

\end{itemize}

\end{discussion}

%%%%%%%%%%%%%%%%%%%%%%%%%%%%%%%%%%%%%%%%%%%%%%%%%%%%%%%%%%%%%%%%%%%%%%

\subsection{2. Métadonnées}

\begin{discussion}

\begin{itemize}

\item Quelles métadonnées vous semblent indispensables pour garantir la
qualité scientifique du corpus ?

\item Certaines informations devraient-elles être évitées afin de limiter
les biais ou les problèmes liés à la confidentialité ?

\item Existe-t-il un format standard que vous recommanderiez pour
l'organisation des données ?

\end{itemize}

\end{discussion}

%%%%%%%%%%%%%%%%%%%%%%%%%%%%%%%%%%%%%%%%%%%%%%%%%%%%%%%%%%%%%%%%%%%%%%

\subsection{3. Prétraitement des données}

\begin{discussion}

\begin{itemize}

\item Quels formats audio recommandez-vous pour ce type de projet
(16 kHz, 44,1 kHz, PCM, FLAC, WAV...) ?

\item Quels contrôles qualité devraient être réalisés avant toute phase
d'apprentissage automatique ?

\item Quels outils ou bibliothèques vous semblent les plus adaptés pour
cette étape ?

\end{itemize}

\end{discussion}

%%%%%%%%%%%%%%%%%%%%%%%%%%%%%%%%%%%%%%%%%%%%%%%%%%%%%%%%%%%%%%%%%%%%%%

\subsection{4. Modélisation}

\begin{discussion}

\begin{itemize}

\item Quels modèles de représentation de la parole vous sembleraient les
plus adaptés pour un corpus de taille limitée ?

\item Recommanderiez-vous une approche auto-supervisée, supervisée ou
hybride dans les premières expérimentations ?

\item Existe-t-il des travaux ou des modèles récents que vous me
conseilleriez d'étudier avant de commencer cette phase ?

\end{itemize}

\end{discussion}

%%%%%%%%%%%%%%%%%%%%%%%%%%%%%%%%%%%%%%%%%%%%%%%%%%%%%%%%%%%%%%%%%%%%%%

\subsection{5. Vision scientifique}

\begin{discussion}

\begin{itemize}

\item La progression proposée dans ALORA vous paraît-elle cohérente du
point de vue scientifique ?

\item Voyez-vous des risques méthodologiques importants dès cette
première phase ?

\item Certaines étapes devraient-elles être réorganisées ou approfondies
avant le début du développement ?

\item Pensez-vous que cette problématique puisse conduire, à terme, à un
véritable sujet de recherche en traitement automatique de la parole ou
en linguistique computationnelle ?

\end{itemize}

\end{discussion}

%%%%%%%%%%%%%%%%%%%%%%%%%%%%%%%%%%%%%%%%%%%%%%%%%%%%%%%%%%%%%%%%%%%%%%

\subsection{Conclusion}

Je suis pleinement conscient que les hypothèses présentées dans ce
document pourront évoluer au cours du projet.

Mon objectif n'est pas de défendre une solution déjà arrêtée, mais de
construire progressivement une méthodologie rigoureuse en m'appuyant sur
les recommandations d'un enseignant-chercheur et sur les résultats
obtenus au fil des expérimentations.

Je serais très heureux d'avoir l'occasion d'échanger avec vous afin de
recueillir vos remarques, vos critiques ainsi que vos suggestions avant
de commencer cette première phase de développement.

%%%%%%%%%%%%%%%%%%%%%%%%%%%%%%%%%%%%%%%%%%%%%%%%%%%%%%%%%%%%%%%%%%%%%%
\section{Feuille de route prévisionnelle}
%%%%%%%%%%%%%%%%%%%%%%%%%%%%%%%%%%%%%%%%%%%%%%%%%%%%%%%%%%%%%%%%%%%%%%

Le projet ALORA est envisagé comme un projet de recherche progressif.
L'objectif n'est pas de développer immédiatement un système complet, mais
de construire, étape par étape, les fondations nécessaires au
développement d'outils fiables pour les langues orales peu dotées.

La feuille de route présentée ci-dessous constitue une première
planification susceptible d'évoluer en fonction des résultats obtenus et
des recommandations du laboratoire.

%%%%%%%%%%%%%%%%%%%%%%%%%%%%%%%%%%%%%%%%%%%%%%%%%%%%%%%%%%%%%%%%%%%%%%

\subsection{Phase 1 : Été 2026}

Cette première phase est consacrée à la construction des fondations du
projet.

Les principaux objectifs sont les suivants :

\begin{itemize}

\item réaliser une étude bibliographique sur les langues peu dotées et
les méthodes modernes de traitement de la parole ;

\item concevoir l'architecture générale du projet ALORA ;

\item développer une première plateforme de collecte collaborative ;

\item définir la structure du corpus et les métadonnées associées ;

\item mettre en place un pipeline de validation et de préparation des
données ;

\item constituer une première version du corpus oral.

\end{itemize}

Les livrables attendus sont :

\begin{itemize}

\item une plateforme fonctionnelle ;

\item un premier corpus structuré ;

\item une documentation technique ;

\item une méthodologie de collecte reproductible.

\end{itemize}

%%%%%%%%%%%%%%%%%%%%%%%%%%%%%%%%%%%%%%%%%%%%%%%%%%%%%%%%%%%%%%%%%%%%%%

\subsection{Phase 2 : Semestre S5}

Cette phase sera consacrée aux premières expérimentations en
apprentissage automatique.

Les travaux envisagés comprennent notamment :

\begin{itemize}

\item l'analyse exploratoire du corpus ;

\item le prétraitement avancé des données audio ;

\item l'étude de modèles de représentation de la parole ;

\item les premières expérimentations sur la découverte automatique de
structures phonologiques ;

\item l'évaluation des approches retenues.

\end{itemize}

L'objectif est de disposer, à la fin du semestre, d'une première base
expérimentale permettant d'orienter les développements futurs.

%%%%%%%%%%%%%%%%%%%%%%%%%%%%%%%%%%%%%%%%%%%%%%%%%%%%%%%%%%%%%%%%%%%%%%

\subsection{Phase 3 : Projet de Fin d'Études}

Le Projet de Fin d'Études constituera la phase principale de
développement scientifique.

Selon les résultats obtenus durant les étapes précédentes, plusieurs
axes pourront être approfondis :

\begin{itemize}

\item amélioration des représentations de la parole ;

\item découverte automatique des unités phonologiques ;

\item apprentissage progressif de la structure de la langue ;

\item étude de méthodes d'induction d'un système d'écriture ;

\item développement de premiers démonstrateurs (transcription,
recherche vocale, traduction ou autres applications).

\end{itemize}

Cette phase permettra également de consolider les résultats obtenus dans
une perspective de publication ou de poursuite en doctorat.

%%%%%%%%%%%%%%%%%%%%%%%%%%%%%%%%%%%%%%%%%%%%%%%%%%%%%%%%%%%%%%%%%%%%%%

\subsection{Perspective à long terme}

Au-delà du cas du soussou, l'ambition d'ALORA est de proposer une
méthodologie générale pouvant être adaptée à d'autres langues orales peu
dotées.

À terme, cette approche pourrait contribuer à la création de nouvelles
ressources linguistiques et faciliter le développement d'outils
numériques pour des langues aujourd'hui très peu représentées dans les
technologies de traitement automatique de la parole.

%%%%%%%%%%%%%%%%%%%%%%%%%%%%%%%%%%%%%%%%%%%%%%%%%%%%%%%%%%%%%%%%%%%%%%

\begin{objectif}

Cette feuille de route ne constitue pas un planning figé.

Elle représente une première proposition de travail destinée à servir de
base de discussion avec le laboratoire. Les différentes étapes pourront
être adaptées en fonction des recommandations scientifiques, des
contraintes rencontrées et des résultats obtenus au cours du projet.

\end{objectif}

%%%%%%%%%%%%%%%%%%%%%%%%%%%%%%%%%%%%%%%%%%%%%%%%%%%%%%%%%%%%%%%%%%%%%%

\section*{Conclusion}

\addcontentsline{toc}{section}{Conclusion}

Le projet ALORA est né d'une réflexion sur les difficultés rencontrées
par les langues orales peu dotées dans le contexte actuel de
l'intelligence artificielle.

Cette note de cadrage présente une première vision du projet ainsi que
les travaux envisagés pour sa phase initiale.

Je suis conscient que plusieurs choix méthodologiques restent ouverts et
que certaines hypothèses pourront évoluer au cours des expérimentations.
C'est précisément pour cette raison que je souhaite bénéficier d'un
encadrement scientifique dès les premières étapes du projet.

J'espère que cette proposition pourra constituer le point de départ
d'une collaboration enrichissante et je serais très heureux d'échanger
avec vous afin de recueillir vos remarques et vos conseils pour faire
évoluer cette idée vers un véritable projet de recherche.

%%%%%%%%%%%%%%%%%%%%%%%%%%%%%%%%%%%%%%%%%%%%%%%%%%%%%%%%%%%%%%%%%%%%%%
%% Bibliographie
%%%%%%%%%%%%%%%%%%%%%%%%%%%%%%%%%%%%%%%%%%%%%%%%%%%%%%%%%%%%%%%%%%%%%%

\newpage

\nocite{*}
\printbibliography

%%%%%%%%%%%%%%%%%%%%%%%%%%%%%%%%%%%%%%%%%%%%%%%%%%%%%%%%%%%%%%%%%%%%%%

\end{document}